\documentclass[book.tex]{subfiles}

\begin{document}

\chapter{The Wave Equation}

\label{chap:wave_equation}
\noindent\rule{11cm}{0.4pt} \\
\textbf{Prerequisites:} \autoref{chap:pde}, \autoref{chap:simple_physics} \\
\textbf{Difficulty Level:} ** \\
\noindent\rule{11cm}{0.4pt}
\vspace{5mm}
The wave equation is the fundamental mathematical model for travelling waves. The equation can be used to model physical waves on a rope, electromagnetic waves, and even solutions to quantum field theories.

\begin{align}
    \frac{\partial^2 u}{\partial t^2} &= c^2 \left ( \frac{\partial^2 u}{\partial x_1^2} + \frac{\partial^2 u}{\partial x_2^2} + \dotsc + \frac{\partial^2 u}{\partial x_n^2}  \right )
\end{align}
We can simplify this expression to
\begin{align}
\frac{\partial^2 u}{\partial t^2} &= c^2 \nabla^2 u
\end{align}
The wave equation is our first example of a hyperbolic differential equation. Hyperbolic equations have a propagation speed (here $c$) that their solutions must obey.

\section{Solving the Wave Equation}

We perform a change of variable to
\begin{align}
\xi &= x - ct \\
\eta &= x + ct
\end{align}

Under this change of variable, the wave equation transforms to
\begin{align}
\frac{\partial^2 u}{\partial \xi \partial \eta}
\end{align}
This equation follows the separated form we studied earlier, so we can write the general solution as
\begin{align}
u(\xi, \eta) &= F(\xi) + G(\eta)
\end{align}
where $F$ and $G$ are arbitrary functions. Changing variables back we can write
\begin{align}
u(x, t) &= F(x-ct) + G(y+ct)
\end{align}
$F$ can be viewed as a right-traveling wave and $G$ as a left-traveling wave. 
\end{document}